\Askhsh{24569_SOLUTION}{1}
\begin{ylh}{1}
    \item [Ύλη από εικόνα]
\end{ylh}
ΛΥΣΗ

\begin{alist}
    \item[α)] Η συνάρτηση $f$ ορίζεται, όταν και μόνο όταν
    \[
        \begin{cases}
            1 - x \ge 0 \\
            \text{και} \\
            1 - \sqrt{1 - x} \ge 0
        \end{cases}
    \]
    Είναι λοιπόν
    \[
        1 - x \ge 0 \Leftrightarrow x \le 1
    \]
    \[
        1 - \sqrt{1 - x} \ge 0 \Leftrightarrow \sqrt{1 - x} \le 1 \Leftrightarrow 1 - x \le 1 \Leftrightarrow x \ge 0
    \]
    Επομένως $0 \le x \le 1$ και ως εκ τούτου $D_f = [0,1]$.

    \item[β)]
    \begin{rlist}
        \item Έστω $x_1,x_2 \in [0,1]$ με $f(x_1) = f(x_2)$. Θα δείξουμε ότι $x_1 = x_2$. Πράγματι έχουμε διαδοχικά:
        \begin{align*}
            f(x_1) &= f(x_2) \\
            \sqrt{1-\sqrt{1-x_1}} &= \sqrt{1-\sqrt{1-x_2}} \\
            1-\sqrt{1-x_1} &= 1-\sqrt{1-x_2} \\
            \sqrt{1-x_1} &= \sqrt{1-x_2} \\
            1-x_1 &= 1-x_2 \\
            x_1 &= x_2 \text{.}
        \end{align*}
        \item Παρατηρούμε ότι $f(0) = 0$. Στη συνέχεια έχουμε:
        \[
            f(f(x)) = 0 \Leftrightarrow f(f(x)) = f(0) \overset{f \text{ "1-1"}}{\Longleftrightarrow} f(x) = 0 \Leftrightarrow f(x) = f(0) \overset{f \text{ "1-1"}}{\Longleftrightarrow} x = 0 \text{.}
        \]
        Επομένως η εξίσωση $f(f(x)) = 0, x \in [0,1]$ έχει μοναδική ρίζα το $0$.
    \end{rlist}
\end{alist}